\chapter{Kurva Bidang}\label{ch:b1:curves}

Kalkulus \cref{ch:b1:derivative,ch:b1:taylor} dibangun bagi
fungsi $y = f(x)$; sedangkan kebanyakan kurva geometri dan mekanika ---
lintasan, lingkaran yang digelindingkan pada lingkaran, orbit --- menolak bentuk itu
dan hadir sebagai kurva yang \emph{terparameter} $t \mapsto (x(t),
y(t))$ atau sebagai kurva \emph{kutub} $r = r(\theta)$. Bab yang klasik
ini mempelajari keduanya, lalu menutup dengan \omterm{def:b1:curves:conics}{irisan kerucut}.

\section{Kurva terparameter}

\begin{definition}\label{def:b1:curves:def}
Sebuah \emph{kurva terparameter}\index{kurva terparameter} adalah \omterm{def:b1:logic:map}{pemetaan} $f
\colon I \to \R^2$, $t \mapsto M(t) = (x(t), y(t))$, dengan $x, y$
berkelas $C^1$ (setidaknya) pada \omterm{prop:b1:reals:intervals}{selang} $I$. Adapun \emph{vektor
kecepatannya} adalah $f'(t) = (x'(t), y'(t))$; titik $M(t_0)$ disebut
\emph{reguler} ketika $f'(t_0) \neq (0,0)$, dan \emph{garis singgungnya}
di situ adalah garis lewat $M(t_0)$ yang berarah $f'(t_0)$.
\end{definition}

\begin{proof}[Mengapa $f'$ mengarahkan garis singgungnya]
Menurut Taylor--Young komponen demi komponen, $M(t) = M(t_0) + (t - t_0) f'(t_0) +
o(t - t_0)$: sehingga arah tali busurnya $\frac{M(t) - M(t_0)}{t - t_0}$
menuju $f'(t_0)$. Adapun hipotesis keregulerannyalah yang membuat
limitnya sebuah \emph{arah}: karena jika $f'(t_0) = 0$, tampilannya runtuh
menjadi $M(t) = M(t_0) + o(t - t_0)$ dan tak mengatakan apa pun tentang bagaimana
kurvanya meninggalkan titik itu --- dan dari situlah perlakuan terpisah atas
titik singular pada metode di bawah, tempat \omterm{def:b1:derivative:def}{turunan} pertama
yang \emph{tak lenyap} mengambil alih peran $f'$.
Perhatikanlah pula bahwa garis singgungnya sebuah objek geometri: karena sebarang
pemarameteran ulang mengubah $f'$ oleh faktor skalar taknol dan
membiarkan garisnya tak berubah.
\end{proof}

\begin{method}[Mempelajari sebuah kurva terparameter]\label{met:b1:curves:study}
\begin{enumerate}
  \item \emph{Reduksikanlah daerah asalnya} lewat kesetangkupan: yaitu hubungan antara
        $M(-t)$, $M(t + T)$, \dots\ dan $M(t)$ (pencerminan,
        translasi, keperiodikan), lalu gambarlah hanya bagian yang tereduksi.
  \item \emph{Ragamnya}: tabelkanlah tanda $x'$ dan $y'$ bersama-sama;
        karena kurvanya bergerak ke kanan/kiri menurut $x'$, dan naik/turun menurut $y'$.
  \item \emph{Titik yang istimewa}: garis singgung mendatar ($y' = 0 \neq
        x'$), garis singgung tegak ($x' = 0 \neq y'$); sedangkan pada titik yang
        \emph{singular} ($f' = 0$), periksalah \omterm{def:b1:derivative:def}{turunan} yang lebih tinggi
        bagi arah singgungnya.
  \item \emph{Perilaku asimtotiknya} pada ujung $I$, lalu sketsalah.
\end{enumerate}
\end{method}

\begin{example}[Cabang asimtotik, dikerjakan]
\label{ex:b1:curves:asymptotes}
$x(t) = t$, $y(t) = t + \dfrac1t$ pada $\intoo{0}{+\infty}$. Ketika $t
\to 0^{+}$: $x \to 0$ sedangkan $y \to +\infty$ --- jadi kurvanya memanjat
sepanjang \emph{asimtot tegak} $x = 0$ (yaitu limit hingga bagi
satu \omterm{prop:b1:vspaces:coordinates}{koordinat}, dan takhingga bagi yang lain). Ketika $t \to +\infty$:
kedua \omterm{prop:b1:vspaces:coordinates}{koordinatnya} meledak, jadi ujilah sebuah garis: karena selisih
\[
y(t) - x(t) = \frac1t \longrightarrow 0^{+}
\]
memperlihatkan \emph{asimtot miring} $y = x$, yang didekati dari
atas. Di antaranya, $y' = 1 - \frac1{t^2}$ lenyap di $t = 1$:
sehingga titik $(1, 2)$ merupakan titik terendah cabangnya, dan $x' = 1 >
0$ sepanjangnya, jadi kurvanya selalu maju ke kanan. Satu busur,
dua asimtot, satu minimum: yaitu gambaran lengkap dari tiga
perhitungan --- dengan resep umumnya kasatmata di bawahnya:
bahwa ketika $x, y \to \infty$ bersama, periksalah $y/x$ bagi calon
kemiringannya (di sini $\to 1$), lalu $y - (\text{kemiringan})\,x$ bagi
perpotongan dan sisi pendekatannya.
\end{example}

\begin{example}[Astroid]\label{ex:b1:curves:astroid}
$x(t) = \cos^3 t$, $y(t) = \sin^3 t$.\index{astroid} Kesetangkupannya:
$M(t + 2\pi) = M(t)$ (jadi pelajarilah satu periode); $M(-t)$ merupakan pencerminan
$M(t)$ pada sumbu $x$; $M(\pi - t)$ pada sumbu $y$; dan $M(\frac\pi2 -
t)$ pada diagonal $y = x$: sehingga cukuplah mempelajari $t \in
\intcc{0}{\frac\pi4}$ lalu membukanya.

Kecepatannya: $f'(t) = 3\sin t\cos t\,(-\cos t, \sin t)$. Pada
$\intoo{0}{\frac\pi2}$ semua titiknya reguler dengan garis singgung yang berarah
$(-\cos t, \sin t)$; sedangkan di $t = 0$ (yaitu titik $(1,0)$) kecepatannya
lenyap: jadi sebuah \emph{titik runcing}\index{titik runcing}, tempat kurvanya berbalik sepanjang
arah singgungnya $(-1, 0)\cdot(\pm)$ --- dan lewat kesetangkupannya keempat
\omterm{ex:b1:curves:astroid}{titik runcingnya} duduk di $(\pm1, 0), (0, \pm1)$.
\end{example}

\begin{example}[Sebuah angka delapan, dipelajari selengkapnya]
\label{ex:b1:curves:lissajous}
$x(t) = \sin t$, $y(t) = \sin 2t$ (yaitu \emph{kurva Lissajous}).
\emph{Kesetangkupannya}: $M(t + \pi) = (-x(t), y(t))$ (yaitu pencerminan pada
sumbu $y$), $M(-t) = (-x(t), -y(t))$ (kesetangkupan pusat),
$M(\pi - t) = (x(t), -y(t))$ (pencerminan pada sumbu $x$): sehingga
cukuplah mempelajari $t \in \intcc{0}{\frac\pi2}$ lalu membukanya.
\emph{Ragamnya}: $x' = \cos t \geq 0$ sepanjangnya, sedangkan $y' =
2\cos 2t$ positif sebelum $t = \frac\pi4$ dan negatif sesudahnya:
jadi busurnya memanjat ke kanan menuju puncak
$\bigl(\frac{\sqrt2}2,\ 1\bigr)$ di $t = \frac\pi4$ (dengan garis singgung
mendatar), lalu menurun ke kanan menuju $(1, 0)$ di $t =
\frac\pi2$, tempat $x' = 0 \neq y'$: jadi garis singgung tegak.
\emph{Titik ganda}: $M(0) = M(\pi) = (0,0)$, dengan dua kecepatan
yang berbeda
\[
f'(0) = (1,\ 2), \qquad f'(\pi) = (-1,\ 2) :
\]
yaitu dua cabang yang \emph{reguler} dan bersilangan di titik asalnya pada sudut
yang berbeda --- jadi titik ganda, bukan titik singular: karena setiap lintasannya
mulus sempurna, dan kedua lintasannya hanya berbagi tempat.
Adapun seluruh kurvanya adalah angka delapan di bawah.
\end{example}

\begin{omfigure}
\begin{tikzpicture}[scale=1.5]
  \draw[->, thin] (-1.25,0) -- (1.3,0) node[below] {$x$};
  \draw[->, thin] (0,-1.25) -- (0,1.3) node[left] {$y$};
  \draw[omDef, very thick, domain=0:360, samples=200, smooth]
    plot ({sin(\x)}, {sin(2*\x)});
  \fill (0.7071, 1) circle (0.035);
  \node[above, font=\small] at (0.7071, 1)
    {$\bigl(\tfrac{\sqrt2}2, 1\bigr)$};
  \fill (0,0) circle (0.035);
\end{tikzpicture}

{\small Kurva Lissajous $(\sin t, \sin 2t)$: dengan sebuah titik ganda
di titik asalnya, tempat dua cabang reguler bersilangan berkecepatan
$(1, 2)$ dan $(-1, 2)$, dan garis singgung mendatar pada keempat
puncaknya. Kesetangkupannya mereduksi seluruh kerjanya menjadi seperempat periode.}
\end{omfigure}

\begin{omfigure}
\begin{tikzpicture}[scale=1.35]
  \draw[->, thin] (-1.3,0) -- (1.35,0) node[below] {$x$};
  \draw[->, thin] (0,-1.3) -- (0,1.35) node[left] {$y$};
  \draw[omDef, very thick, domain=0:360, samples=200, smooth]
    plot ({cos(\x)^3}, {sin(\x)^3});
  \node[below right, font=\small] at (1,0) {$1$};
  \node[above left, font=\small] at (0,1) {$1$};
\end{tikzpicture}

{\small Astroid $(\cos^3 t, \sin^3 t)$: yaitu empat busur yang bertemu pada
empat \omterm{ex:b1:curves:astroid}{titik runcing}. Ialah kurva yang dijejaki sebuah titik pada lingkaran berjari-jari
$\frac14$ yang menggelinding di dalam lingkaran satuan.}
\end{omfigure}

\section{Kurva kutub}

\begin{definition}\label{def:b1:curves:polar}
Sebuah \emph{kurva kutub}\index{kurva kutub} diberikan oleh $r = r(\theta)$:
dengan titik berparameter $\theta$ adalah
\[
M(\theta) = r(\theta)\,\vec u(\theta),
\qquad \vec u(\theta) = (\cos\theta, \sin\theta).
\]
Dengan $\vec v(\theta) = (-\sin\theta, \cos\theta) = \vec
u\,'(\theta)$, kecepatannya adalah
\[
M'(\theta) = r'(\theta)\, \vec u(\theta) + r(\theta)\, \vec
v(\theta) .
\]
Akibatnya: di tempat $r \neq 0$, titiknya reguler, dan
garis singgungnya bersudut $V$ dengan sinar garisnya, yang diberikan oleh $\tan V =
\frac{r}{r'}$ (yaitu sudut antara $M'$ dan $\vec u$); sedangkan di tempat $r(\theta_0)
= 0$, kurvanya lewat titik asalnya \emph{dengan garis singgungnya sinar garis
$\theta = \theta_0$} (yaitu berarah $\vec u(\theta_0)$, yang terbaca dari
$M'(\theta_0) = r'(\theta_0)\vec u(\theta_0)$, atau dari tali busur
limitnya: karena tali busur dari $O$ ke $M(\theta)$ dibawa oleh $\vec
u(\theta)$ sendiri, yang menuju $\vec u(\theta_0)$ --- sehingga
kaidahnya berlaku bahkan ketika $r'(\theta_0) = 0$ dan titiknya
singular, dan itulah sebabnya \omterm{ex:b1:curves:astroid}{titik runcing} kutub di titik asalnya, seperti pada
kardioid, memperoleh garis singgungnya secara cuma-cuma).
\end{definition}

\begin{example}[Kardioid]\label{ex:b1:curves:cardioid}
$r(\theta) = 1 + \cos\theta$.\index{kardioid} Kesetangkupannya: $r(-\theta)
= r(\theta)$: yaitu pencerminan pada sumbu $x$; jadi pelajarilah $\theta \in
\intcc{0}{\pi}$. Adapun $r$ menurun dari $2$ ke $0$; dan di $\theta = \pi$,
$r = 0$: sehingga kurvanya mencapai titik asalnya secara menyinggung sinar garis
$\theta = \pi$ (yaitu sumbu $x$), lalu membentuk \omterm{ex:b1:curves:astroid}{titik runcing} di situ --- yakni ujung
hatinya. Sedangkan garis singgung di $\theta = 0$: $r' = 0$, sehingga $\tan V = \infty$:
jadi tegak lurus sumbunya.

Adapun sudut $V$ pantas dibaca sekali lagi. Di $\theta =
\frac\pi2$: $r = 1$ dan $r' = -1$, sehingga $\tan V = \frac{r}{r'} =
-1$: jadi garis singgungnya bersudut tiga perempat sudut siku dengan
sinar garis yang keluar --- karena kurvanya sudah membengkok kembali menuju
\omterm{ex:b1:curves:astroid}{titik runcingnya}. Rumus $\tan V = r/r'$ menyampaikan arah singgungnya
sepanjang seluruh kurvanya tanpa perhitungan $M'(\theta)$
sama sekali: jadi ia padanan kutub bagi membaca sebuah kemiringan.
\end{example}

\begin{example}[Kutub ke kartesius: sebuah lingkaran yang tersembunyi]
\label{ex:b1:curves:polarcircle}
Apakah \omterm{def:b1:curves:polar}{kurva kutub} $r = 2\cos\theta$ itu? Kalikanlah dengan $r$:
$r^2 = 2r\cos\theta$, yaitu $x^2 + y^2 = 2x$, yakni
\[
(x - 1)^2 + y^2 = 1 :
\]
yaitu lingkaran berpusat $(1, 0)$ dan berjari-jari $1$, yang lewat
titik asalnya. Pembukuannya penting: karena ketika $\theta$ menjelajahi
$\intoc{-\frac\pi2}{\frac\pi2}$ seluruh lingkarannya tersapu
tepat sekali ($r$ lenyap pada kedua ujungnya), dan di $\theta =
\pm\frac\pi2$ kaidah ``garis singgung di titik asalnya sepanjang sinar garis
$\theta = \theta_0$'' memberikan garis singgung yang \emph{tegak} di situ ---
yang cocok dengan geometrinya, karena sumbu tegaknya memang menyinggung
lingkaran ini di $O$. Untuk $\theta$ di luar jangkauan itu, $r < 0$
menjejaki ulang lingkaran yang sama: yaitu pengingat bahwa \omterm{def:b1:curves:polar}{kurva kutub} merupakan
objek yang \emph{terparameter}, yang boleh melintasi dirinya sendiri.
\end{example}

\begin{omfigure}
\begin{tikzpicture}[scale=1.5]
  \draw[->, thin] (-1.2,0) -- (1.3,0) node[below] {$x$};
  \draw[->, thin] (0,-1.2) -- (0,1.3) node[left] {$y$};
  \draw[omProp, very thick, domain=0:360, samples=300, smooth]
    plot ({cos(2*\x)*cos(\x)}, {cos(2*\x)*sin(\x)});
  \draw[thin, dashed] (-0.85,-0.85) -- (0.85,0.85);
  \draw[thin, dashed] (-0.85,0.85) -- (0.85,-0.85);
\end{tikzpicture}

{\small Mawar berkelopak empat $r = \cos 2\theta$
(\cref{exo:b1:curves:4}). Adapun kelopak sepanjang sumbu $y$
terjejaki dengan $r < 0$ (karena titiknya terplot pada sinar garis yang berlawanan dengan
$\theta$); sedangkan diagonal putus-putus $\theta = \pm\frac\pi4$ merupakan
garis singgung di titik asalnya, tempat $r$ lenyap.}
\end{omfigure}

\begin{omfigure}
\begin{tikzpicture}[scale=1.15]
  \draw[->, thin] (-0.7,0) -- (2.4,0) node[below] {$x$};
  \draw[->, thin] (0,-1.7) -- (0,1.7) node[left] {$y$};
  \draw[omThm, very thick, domain=0:360, samples=200, smooth]
    plot ({(1+cos(\x))*cos(\x)}, {(1+cos(\x))*sin(\x)});
  \node[below, font=\small] at (2,0) {$2$};
\end{tikzpicture}

{\small Kardioid $r = 1 + \cos\theta$. \omterm{def:b1:curves:polar}{Kurva kutub} dibaca dengan
menyapu sudutnya: karena jari-jarinya membengkak dan menyusut ketika $\theta$
berputar.}
\end{omfigure}

\section{Panjang busur}

\begin{definition}[Panjang sebuah busur]\label{def:b1:curves:arclength}
Adapun \emph{panjang}\index{panjang busur} sebuah busur $C^1$ yaitu $f \colon
\intcc{a}{b} \to \R^2$ adalah \omterm{thm:b1:integration:def}{integral} atas lajunya:
\[
L = \int_a^b \norm{f'(t)}\,\dd t
= \int_a^b \sqrt{x'(t)^2 + y'(t)^2}\;\dd t .
\]
(Alasannya: pada \omterm{prop:b1:reals:intervals}{selang} yang kecil, $M(t + h) \approx M(t) + h
f'(t)$, sehingga busurnya dekat pada segi banyak yang panjang ruasnya berjumlah
menjadi sebuah \omterm{thm:b1:integration:riemann}{jumlah Riemann} atas $\norm{f'}$,
\cref{thm:b1:integration:riemann}.) Adapun bagi \omterm{def:b1:curves:polar}{kurva kutub} $r =
r(\theta)$, kecepatannya $r'\vec u + r\vec v$ berkomponen
\omterm{def:b1:euclid:orthogonal}{ortogonal}, sehingga
\[
L = \int_{\theta_1}^{\theta_2}
\sqrt{r'(\theta)^2 + r(\theta)^2}\;\dd\theta .
\]
Panjangnya tak bergantung pada pemarameteran (yang monoton dan $C^1$)
yang dipilih: karena menyulihkan $t = \varphi(s)$ ke dalam
\omterm{thm:b1:integration:def}{integralnya} (\cref{thm:b1:integration:parts}) mengalikan $f'$ dengan
$\varphi'$ dan $\dd t$ dengan $\varphi'^{-1}$.
\end{definition}

\begin{example}[Periksa kewarasan: lingkarannya]\label{ex:b1:curves:circlelength}
Untuk $f(t) = (R\cos t, R\sin t)$ pada $\intcc{0}{2\pi}$: $\norm{f'} =
R$, sehingga $L = 2\pi R$ --- jadi definisinya mengembalikan kelilingnya.
Uji pemarameteran ulangnya: \omterm{def:b1:logic:map}{pemetaan} $g(t) = (R\cos 2t, R\sin 2t)$ pada
$\intcc{0}{\pi}$ menggambar lingkaran yang \emph{sama} pada laju yang berganda
$\norm{g'} = 2R$, dan
\[
\int_0^{\pi} 2R\,\dd t = 2\pi R
\]
lagi: separuh waktunya, dua kali lajunya, \omterm{def:b1:curves:arclength}{panjang} yang sama --- yaitu
keinvarianan yang dijanjikan pada definisinya, disaksikan sekali pada bilangan.
(Menjalankan $g$ pada seluruh $\intcc{0}{2\pi}$ akan memberikan $4\pi R$: karena
kurva yang dilintasi dua kali dua kali lebih \omterm{def:b1:curves:arclength}{panjang} \emph{sebagai perjalanan};
sebab \omterm{def:b1:curves:arclength}{panjangnya} mengukur lintasan yang terparameter, dan pembukuan yang jujur atas
\omterm{prop:b1:reals:intervals}{selangnya} bagian dari perhitungannya.)
Bagi astroid $(\cos^3t, \sin^3t)$: $\norm{f'} = 3\abs{\sin
t\cos t} = \tfrac32\abs{\sin 2t}$, dan lewat kesetangkupannya $L = 4\int_0^{\pi/2}
\tfrac32\sin 2t\,\dd t = 6$: yaitu kurva yang digambar di dalam lingkaran satuan,
yang berpanjang $6 < 2\pi$. Adapun soal akhir pekannya mengukur lengkungan
yang paling termasyhur.
\end{example}

\begin{omfigure}
\begin{tikzpicture}[scale=0.85]
  \draw[->, thin] (-0.5,0) -- (7.2,0) node[below] {$x$};
  \draw[->, thin] (0,-0.4) -- (0,2.6) node[left] {$y$};
  \draw[omThm, very thick, domain=0:6.2832, samples=160, smooth]
    plot ({\x - sin(\x r)}, {1 - cos(\x r)});
  \draw[omDef, thick] (2,1) circle (1);
  \fill[omDef] (2,1) circle (0.04);
  \fill (2 - 0.9093, 1.4161) circle (0.06);
  \draw[omProp, thick] (2,0) -- (2 - 0.9093, 1.4161) -- (2,2);
  \fill (2,0) circle (0.05);
  \fill (2,2) circle (0.05);
  \node[below, font=\small] at (2,0) {$C$};
  \node[above, font=\small] at (2,2) {$T$};
  \node[left, font=\small] at (2 - 0.9093, 1.4161) {$M$};
  \node[below, font=\small] at (6.2832,0) {$2\pi R$};
\end{tikzpicture}

{\small Satu lengkungan sikloid (berjari-jari $R = 1$), lingkaran
yang menggelinding di $t = 2$, dan kedua garis pemandu soal akhir pekannya:
karena tali busur $MC$ ke titik sentuhnya bersifat \emph{normal} terhadap
kurvanya, sedangkan tali busur $MT$ ke puncak lingkarannya bersifat \emph{menyinggung}.}
\end{omfigure}

\section{Irisan kerucut}

\begin{definition}[Definisi fokus--direktriks]\label{def:b1:curves:conics}
Tetapkanlah sebuah titik $F$ (\emph{fokus}), sebuah garis $D$ yang tak lewat $F$
(\emph{direktriks}) dan $e > 0$ (\emph{eksentrisitas}). Maka
\emph{irisan kerucut}\index{irisan kerucut} atas data ini adalah
\[
\mathcal{C} = \{M : d(M, F) = e\; d(M, D)\}:
\]
yaitu \emph{elips} untuk $e < 1$, \emph{parabola} untuk $e = 1$, dan
\emph{hiperbola} untuk $e > 1$. (Adapun lingkarannya muncul sebagai limit
yang merosot $e \to 0$.)
\end{definition}

\begin{example}[Definisinya, diperiksa pada sebuah parabola]
\label{ex:b1:curves:parabolacheck}
Ambillah parabola $y^2 = 4x$, yaitu $2p = 4$: dengan fokus $F = (1, 0)$
dan direktriks $D : x = -1$ (karena bentuk tereduksinya pada
\cref{thm:b1:curves:reduced} menempatkannya di $\pm\frac p2$). Pada
titik $M = (1, 2)$ kurvanya:
\[
MF = \sqrt{(1-1)^2 + 2^2} = 2,
\qquad
d(M, D) = 1 - (-1) = 2 :
\]
sama, sebagaimana dituntut $e = 1$. Di $M' = (4, 4)$: $MF' = \sqrt{9 +
16} = 5$ dan $d(M', D) = 5$ lagi. Jadi definisi
fokus--direktriksnya bukan abstraksi: melainkan sepasang jarak yang dapat
diukur orang pada sebarang titik, dan aljabar persamaan
tereduksinya tak lain pengukuran ini yang dikerjakan sekali untuk
selamanya.
\end{example}

\begin{theorem}[Persamaan tereduksi]\label{thm:b1:curves:reduced}
Pada kerangka \omterm{def:b1:euclid:orthogonal}{ortonormal} yang terpilih baik:
\[
\text{elips: } \frac{x^2}{a^2} + \frac{y^2}{b^2} = 1
\quad (a \geq b > 0),
\qquad
\text{hiperbola: } \frac{x^2}{a^2} - \frac{y^2}{b^2} = 1,
\qquad
\text{parabola: } y^2 = 2px,
\]
dengan, bagi elipsnya: fokus di $(\pm c, 0)$, $c = \sqrt{a^2 - b^2}$,
$e = \frac ca$, dan pencirian bifokalnya $MF + MF' = 2a$;
sedangkan bagi hiperbolanya: $c = \sqrt{a^2 + b^2}$, $e = \frac ca$,
$\abs{MF - MF'} = 2a$, dan asimtot $y = \pm\frac ba x$.
\end{theorem}

\begin{proof}
Ambillah fokusnya di titik asalnya dan direktriksnya tegak, $x = -h$
($h > 0$): maka syarat $MF^2 = e^2 d(M, D)^2$ terbaca $x^2 + y^2 =
e^2 (x + h)^2$. Untuk $e = 1$: $y^2 = 2hx + h^2$, yaitu parabola setelah
geseran $x \mapsto x - \frac h2$ (sehingga $p = h$). Untuk $e \neq 1$:
dengan melengkapkan kuadratnya pada $x$,
\[
(1 - e^2)\Bigl(x - \frac{e^2 h}{1 - e^2}\Bigr)^{\!2} + y^2
= e^2h^2 + \frac{e^4 h^2}{1 - e^2}
= \frac{e^2 h^2}{1 - e^2} .
\]
Tetapkanlah $X = x - \frac{e^2h}{1-e^2}$ (yaitu geseran kerangkanya). Ketika $e < 1$,
bagilah dengan ruas kanannya yang positif:
$\frac{X^2}{a^2} + \frac{y^2}{b^2} = 1$ dengan $a = \frac{eh}{1 -
e^2}$, $b = \frac{eh}{\sqrt{1-e^2}}$; sedangkan ketika $e > 1$, kedua ruas
tampilannya berbalik tanda sebagaimana mestinya dan pembagian yang sama memberikan
$\frac{X^2}{a^2} - \frac{y^2}{b^2} = 1$ dengan $a = \frac{eh}{e^2 -
1}$, $b = \frac{eh}{\sqrt{e^2 - 1}}$. Adapun nilai $c$ yang dinyatakannya
menyusul ($c^2 = a^2 - b^2$ atau $a^2 + b^2$ memberikan $c = ea$ pada kedua
kasusnya, yang menempatkan fokusnya dengan benar), sedangkan sifat bifokalnya berupa
pemeriksaan langsung pada persamaan tereduksinya. Terperincinya bagi
elipsnya, dengan $F' = (c, 0)$ dan $M = (X, y)$ pada kurvanya:
\[
MF'^2 = (X - c)^2 + y^2
= X^2 - 2cX + c^2 + b^2\Bigl(1 - \frac{X^2}{a^2}\Bigr)
= \frac{c^2}{a^2}X^2 - 2cX + a^2
= (a - eX)^2 ,
\]
dengan memakai $c^2 + b^2 = a^2$ dan $e = \frac ca$; dan karena $\abs X \leq
a$ dan $e < 1$, maka $a - eX > 0$, sehingga $MF' = a - eX$ tanpa akar
kuadrat yang pernah ditarik. Adapun perhitungan cerminnya memberikan $MF = a + eX$,
dan dari situlah $MF + MF' = 2a$, yang tetap. Bagi hiperbolanya aljabar
yang sama menghasilkan $MF = \abs{a + eX}$ dan $MF' = \abs{eX - a}$,
dengan selisih $\pm2a$ menurut cabangnya.
\end{proof}

\begin{example}\label{ex:b1:curves:conicexample}
$\frac{x^2}{25} + \frac{y^2}{9} = 1$: yaitu elips, $a = 5$, $b = 3$, $c
= 4$: dengan fokus $(\pm4, 0)$, dan eksentrisitas $\frac45$. Pemarameterannya:
$(5\cos t, 3\sin t)$ --- yaitu lingkaran yang direntangkan secara anisotropis; dan
jumlah jarak ke fokusnya dari sebarang titiknya adalah $10$.
\end{example}

\begin{example}[Membaca sebuah orbit dari persamaan kutubnya]
\label{ex:b1:curves:orbit}
\omterm{def:b1:curves:conics}{Irisan kerucut} kutub $r = \dfrac{1}{1 + \frac12\cos\theta}$
(\cref{exo:b1:curves:7} dengan $p = 1$, $e = \frac12$) merupakan
elips dengan \emph{fokus di titik asalnya} --- yaitu geometri sebuah
orbit planet dengan mataharinya di $O$. Tariklah segalanya dari
$p$ dan $e$: karena dari reduksi
\cref{thm:b1:curves:reduced}, $a = \dfrac{p}{1 - e^2} =
\dfrac{1}{3/4} = \dfrac43$ dan $c = ea = \dfrac23$. Adapun kedua
apsisnya memeriksanya tanpa teori sama sekali:
\[
r(0) = \frac{1}{3/2} = \frac23 = a - c
\quad (\text{perihelion}),
\qquad
r(\pi) = \frac{1}{1/2} = 2 = a + c
\quad (\text{aphelion}),
\]
dan jumlahnya $\frac23 + 2 = \frac83 = 2a$ memulihkan sumbu
mayornya. Adapun bentuk kutubnya yang alami kapan pun sebuah fokus
terbedakan secara fisik; sedangkan bentuk kartesius tereduksinya, kapan pun
sumbu kesetangkupannya terbedakan. Mengalihkan antara keduanya tepat
apa yang dikerjakan perhitungan kuadrat lengkap pada teoremanya.

Akhirnya, eksentrisitasnya sebagai sebuah tombol: jagalah $p = 1$ lalu putarlah $e$ pada $r
= \frac{1}{1 + e\cos\theta}$. Di $e = 0$: yaitu lingkaran $r = 1$. Di
$e = \frac12$: yaitu elips yang baru dipelajari, dengan $r$ berayun antara
$\frac23$ dan $2$. Di $e = 1$: $r(\theta) \to \infty$ ketika $\theta
\to \pi$ --- sehingga kurvanya tak lagi menutup: jadi sebuah parabola, dengan
titik terjauhnya terdorong ke takhingga. Di $e = 2$: penyebutnya
lenyap di $\cos\theta = -\frac12$, dan hanya $\theta \in
\intoo{-\frac{2\pi}3}{\frac{2\pi}3}$ yang bertahan: yaitu satu cabang
hiperbola, yang lolos sepanjang dua arah asimtotik. Jadi satu
rumus, seluruh keluarga \omterm{def:b1:curves:conics}{irisan kerucutnya}, dan titik peralihannya $e =
1$ kasatmata sebagai saat penyebutnya pertama kali mencapai nol.
\end{example}

\begin{remark}[Jebakan yang lazim]
\label{rem:b1:curves:pitfalls}
\emph{Titik ganda bukanlah titik singular}: karena pada sebuah
persilangan diri, setiap cabangnya reguler; sedangkan ``singular'' menunjuk pada
$f'(t_0) = 0$ bagi satu nilai parameternya
(\cref{ex:b1:curves:lissajous} lawan \omterm{ex:b1:curves:astroid}{titik runcing} astroidnya).
\emph{Garis singgung tegak lawan \omterm{ex:b1:curves:astroid}{titik runcing}}: $x'(t_0) = 0 \neq y'(t_0)$
merupakan titik reguler bergaris singgung tegak; hanya $x' = y' = 0$
yang menuntut penjabaran berorde lebih tinggi.
\emph{Jari-jari negatif terplot pada sinar garis yang berlawanan}: karena untuk $r(\theta) <
0$ titiknya adalah $-\abs{r}\,\vec u(\theta)$, bersudut $\theta +
\pi$ (\cref{exo:b1:curves:5}); dan melupakan ini kehilangan simpul dalamnya
atau menggandakan kurvanya.
\emph{\omterm{def:b1:curves:arclength}{Panjang busur} mengintegralkan $\norm{f'}$, bukan $f'$}: belahlah
\omterm{thm:b1:integration:def}{integralnya} pada nol lajunya --- karena bagi astroidnya,
mengintegralkan $\frac32\sin 2t$ atas satu periode penuh tanpa nilai
mutlak memberikan $0$, bukan $6$.
\emph{\omterm{def:b1:curves:conics}{Irisan kerucut} mesti direduksi sebelum dibaca}: karena pada $9x^2 +
25y^2 - 36x - 50y - 164 = 0$, baik sumbunya, pusatnya, maupun
eksentrisitasnya tak kasatmata sampai kuadratnya dilengkapkan
(\cref{exo:b1:curves:6}); dan bagian kuadratik yang lenyap pada satu
variabel berarti parabola, bukan ``elips yang merosot''.
\end{remark}

\begin{remark}[Ke mana kurva ini pergi]
\label{rem:b1:curves:whereused}
\omterm{def:b1:curves:def}{Kurva terparameter} adalah bahasa mekanika: karena lintasan
berupa kurva, vektor kecepatan berupa vektor singgung, dan \omterm{def:b1:curves:arclength}{panjang
busur} \cref{def:b1:curves:arclength} adalah jarak yang ditempuh. Adapun
\omterm{def:b1:curves:conics}{irisan kerucutnya} muncul kembali di mana pun hukum kuadrat terbalik bekerja --- karena
orbit planet berupa elips dengan mataharinya pada sebuah fokus. Pada
\cref{ch:b1:multivar}, kurvanya menjadi \omterm{def:b1:logic:sets}{himpunan} aras fungsi
dua variabel, dan garis singgung bab ini bertemu
gradien bab berikutnya. Adapun jilid Tahun~2 menambahkan kelengkungan dan
bentuk kanonik lokal sebuah kurva; sedangkan soal akhir pekan di bawah
sudah menarik, hanya dengan perkakas tahun ini, segala yang
diketahui abad ketujuh belas tentang kurvanya yang paling termasyhur.
\end{remark}

\begin{remark}[Cakrawala di dalam Buku 3]
\label{rem:b1:curves:perspectives}
Bab ini melahap seluruh paruh pertama jilid ini dan
menyuapi bab terakhirnya. \emph{Yang dilahap}: vektor singgungnya adalah
\omterm{def:b1:derivative:def}{turunan} (\cref{ch:b1:derivative}), \omterm{def:b1:curves:arclength}{panjang busur} dan luasnya adalah
\omterm{thm:b1:integration:def}{integral} (\cref{ch:b1:integration}), \omterm{ex:b1:curves:astroid}{titik runcingnya} diselesaikan lewat
penjabaran Taylor (\cref{ch:b1:taylor}), tautokronnya adalah
persamaan diferensial \omterm{def:b1:linmaps:def}{linear} (\cref{ch:b1:diffeq}), dan setiap
jarak dan sudutnya bersifat Euklides (\cref{ch:b1:euclid}).
\emph{Yang disuapi}: pada \cref{ch:b1:multivar}, sebuah kurva $t \mapsto (x(t),
y(t))$ yang digambar di dalam \omterm{def:b1:logic:sets}{himpunan} aras $\{f = c\}$ diturunkan lewat
kaidah rantai, dan kesamaan yang dihasilkannya $\langle \nabla f,\
f'(t)\rangle = 0$ menikahkan vektor singgung bab ini dengan
gradien bab berikutnya --- karena garis singgung kurvanya dan
arah normal permukaannya merupakan perhitungan yang sama, yang dilihat
dari kedua tepiannya.
\end{remark}

\section{Latihan}

\begin{exercise}[$\star$]\label{exo:b1:curves:1}
Pelajarilah lalu sketsalah kurva $x(t) = t^2$, $y(t) = t^3$ (kesetangkupan,
ragam, dan perilakunya pada titik singular $t = 0$).
\end{exercise}

\begin{exercise}[$\star$]\label{exo:b1:curves:2}
Bagi sikloid $x(t) = t - \sin t$, $y(t) = 1 - \cos t$ (yaitu
lintasan sebuah titik pada roda yang menggelinding berjari-jari $1$): kenalilah
kesetangkupan translasi periodiknya, titik singularnya, dan
arah singgungnya di $t = 0$ \emph{(jabarkanlah $x$ dan $y$ sampai orde
taknol yang pertama)}.
\end{exercise}

\begin{exercise}[$\star$]\label{exo:b1:curves:3}
Berikanlah garis singgung kurva $(\cos^3 t, \sin^3 t)$ di $t =
\frac\pi4$, lalu periksalah bahwa ruas garis singgung ini yang dipotong
sumbunya berpanjang $1$ --- yaitu sifat astroid yang termasyhur (yang benar pada
setiap titik regulernya).
\end{exercise}

\begin{exercise}[$\star$]\label{exo:b1:curves:4}
Sketsalah \omterm{def:b1:curves:polar}{kurva kutub} $r = \cos 2\theta$ (yaitu mawar berkelopak empat) dan
$r = \frac{1}{\cos\theta}$ pada $\intoo{-\frac\pi2}{\frac\pi2}$
(kenalilah sebuah garis).
\end{exercise}

\begin{exercise}[$\star\star$]\label{exo:b1:curves:5}
Pelajarilah \omterm{def:b1:curves:polar}{kurva kutub} $r = 1 + 2\cos\theta$ (yaitu sebuah limason): daerah
tempat $r \geq 0$ lawan $r < 0$ (karena titik yang diplot dengan jari-jari negatif
duduk pada sinar garis yang berlawanan), lintasannya lewat titik asalnya, simpul dalamnya,
lalu sketsalah.
\end{exercise}

\begin{exercise}[$\star\star$]\label{exo:b1:curves:6}
Kenalilah \omterm{def:b1:curves:conics}{irisan kerucut} $9x^2 + 25y^2 - 36x - 50y - 164 = 0$: reduksikanlah dengan
melengkapkan kuadratnya, lalu berikanlah pusat, setengah sumbu, fokus, dan eksentrisitasnya.
\end{exercise}

\begin{exercise}[$\star\star$]\label{exo:b1:curves:7}
Buktikan bahwa persamaan kutub sebuah \omterm{def:b1:curves:conics}{irisan kerucut} yang berfokus di titik asalnya
adalah
\[
r = \frac{p}{1 + e\cos\theta}
\]
(dengan direktriks tegak pada jarak $\frac pe$ di kanan fokusnya).
Nilai $\theta$ yang mana yang diperkenankan ketika $e > 1$?
\end{exercise}

\begin{exercise}[$\star\star$]\label{exo:b1:curves:8}
Dari sifat bifokal $MF + MF' = 2a$, simpulkanlah ``konstruksi
tukang kebun'' bagi elipsnya, lalu buktikanlah bahwa garis singgung di $M$
bersudut sama dengan $MF$ dan $MF'$ \emph{(yaitu sifat pemantulannya;
turunkanlah $\norm{M(t) - F} + \norm{M(t) - F'} = 2a$ lalu
tafsirkanlah jumlah hasil kali vektor satuannya yang lenyap)}.
\end{exercise}

\begin{exercise}[$\star\star\star$]\label{exo:b1:curves:9}
\emph{Lemniskat Bernoulli} adalah \omterm{def:b1:curves:polar}{kurva kutub} $r^2 = \cos
2\theta$ (ambillah $r = \sqrt{\cos2\theta}$ di tempat ia terdefinisi).
\begin{enumerate}
  \item Berikanlah daerah asal, kesetangkupan, dan garis singgungnya di titik asalnya, lalu
        sketsalah ia.
  \item Buktikan bahwa ia tempat kedudukan titik $M$ dengan $MF \cdot MF'
        = \frac12$ dengan $F, F' = \bigl(\pm\frac{1}{\sqrt2},
        0\bigr)$. \emph{(Hitunglah $MF^2\,MF'^2$ dalam \omterm{prop:b1:vspaces:coordinates}{koordinat}
        kutub.)}
\end{enumerate}
\end{exercise}

\begin{exercise}[$\star\star$]\label{exo:b1:curves:10}
Hitunglah \omterm{def:b1:curves:arclength}{panjang} total kardioid $r = 1 + \cos\theta$
dengan memakai rumus kutub \cref{def:b1:curves:arclength}
\emph{(karena kesamaan $1 + \cos\theta = 2\cos^2\frac\theta2$ mengubah
akar kuadratnya menjadi $2\abs{\cos\frac\theta2}$)}.
\end{exercise}

\begin{exercise}[$\star\star$]\label{exo:b1:curves:11}
Tunjukkan bahwa garis singgung elips $(a\cos t, b\sin t)$ pada
titik berparameter $t$ berpersamaan
\[
\frac{x\cos t}{a} + \frac{y\sin t}{b} = 1 ,
\]
lalu simpulkanlah garis singgung $\frac{x^2}{a^2} + \frac{y^2}{b^2} = 1$
pada sebuah titik $(x_0, y_0)$ elipsnya: $\frac{x\,x_0}{a^2} +
\frac{y\,y_0}{b^2} = 1$ (yaitu kaidah ``pembelahan kuadratnya'').
\end{exercise}

\begin{exercise}[$\star\star\star$]\label{exo:b1:curves:12}
\emph{Spiral logaritmik} adalah \omterm{def:b1:curves:polar}{kurva kutub} $r =
\eu^{k\theta}$ (dengan $k > 0$ yang tetap, $\theta \in \R$).
\begin{enumerate}
  \item Tunjukkan bahwa sudut $V$ antara jari-jarinya dan
        garis singgungnya bersifat \emph{tetap} ($\tan V = \frac1k$) --- sehingga
        spiralnya menyilang setiap sinar garis pada sudut yang sama.
  \item Tunjukkan bahwa memutar spiralnya sebesar sudut $c$ memetakannya
        ke penskalaannya oleh faktor $\eu^{-kc}$: jadi setiap
        rotasi spiralnya merupakan perbesarannya
        (yaitu keserupaan diri).
  \item Hitunglah \omterm{def:b1:curves:arclength}{panjang busur} $\theta \in
        \intoc{-\infty}{\theta_0}$ (sebagai limit \omterm{def:b1:curves:arclength}{panjangnya} pada
        $\intcc{A}{\theta_0}$, $A \to -\infty$) lalu amatilah bahwa
        ia hingga: yaitu kurva yang berpilin takhingga banyaknya kali
        mengelilingi titik asalnya, tetapi berpanjang hingga.
\end{enumerate}
\end{exercise}

\section{Soal: sikloid, ratu segala kurva}

\begin{problem}\label{pb:b1:curves:1}
Sebuah roda berjari-jari $R$ menggelinding tanpa tergelincir sepanjang sumbu $x$;
dan titik pelek yang mula-mula di titik asalnya menjejaki
\emph{sikloid}. Abad ketujuh belas bertengkar atas kurva ini
--- Galileo menimbang guntingan kertasnya, Wren mengukurnya,
Roberval menghitung luasnya, Huygens membangun jam di atasnya --- dan
setiap hasil mereka terjangkau bab ini. Kita
membuktikan keempat klasiknya: yaitu konstruksi singgungnya, \omterm{def:b1:curves:arclength}{panjang} Wren
$8R$, luas $3\pi R^2$, dan sifat tautokron
Huygens.\index{sikloid}\index{tautokron}

\medskip
\noindent\textbf{Bagian I --- Menggelinding dan garis singgungnya.}
\begin{enumerate}
  \item Setelah rodanya berputar sebesar sudut $t$, pusatnya
        duduk di $\Omega(t) = (Rt, R)$ (karena menggelinding tanpa tergelincir:
        jarak sentuhnya $=$ busur yang tergelinding). Tunjukkan bahwa titik
        yang ditandai berada di
        \[
        M(t) = \bigl(R(t - \sin t),\; R(1 - \cos t)\bigr).
        \]
  \item Hitunglah $f'(t)$ lalu tunjukkanlah $\norm{f'(t)} =
        2R\,\abs{\sin\frac t2}$; lalu carilah titik singularnya
        (yaitu \omterm{ex:b1:curves:astroid}{titik runcing} --- bandingkanlah \cref{exo:b1:curves:2}) dan puncak
        setiap lengkungannya.
  \item Misalkan $C(t) = (Rt, 0)$ titik sentuhnya dan $T(t) =
        (Rt, 2R)$ puncak rodanya. Buktikan bahwa untuk $0 < t <
        2\pi$ vektor $M - C$ bersifat \emph{normal} terhadap kurvanya di
        $M(t)$ sedangkan vektor $T - M$ bersifat \emph{menyinggung}: jadi untuk menggambar
        garis singgung sebuah sikloid, hubungkanlah titiknya ke puncak
        lingkaran yang menggelindingkannya. \emph{(Faktorkanlah segalanya lewat
        $\sin\frac t2$ dan $\cos\frac t2$.)}
  \item Tafsirkanlah pertanyaan 3 secara kinematika: bahwa titik sentuhnya adalah
        \emph{pusat rotasi sesaatnya}, dan
        laju $M$ sama dengan jaraknya ke $C$ (bagi kecepatan sudut
        satuan). Periksalah $\norm{M - C} = 2R\abs{\sin\frac t2}$.
  \item Buktikan hubungan tinggi--laju
        \[
        \norm{f'(t)}^2 = 2R\,y(t) :
        \]
        bahwa pada sikloid yang dilintasi pada kecepatan sudut satuan,
        lajunya pada setiap titik tepat laju jatuh \omterm{def:b1:vspaces:free}{bebas} bagi
        penurunan yang sama dengan tinggi saat itu. (Simpanlah ini bagi
        Bagian IV.)
\end{enumerate}

\medskip
\noindent\textbf{Bagian II --- Teorema Wren: lengkungannya berpanjang
$8R$.}
\begin{enumerate}[resume]
  \item Dengan memakai \cref{def:b1:curves:arclength}, hitunglah \omterm{def:b1:curves:arclength}{panjang}
        satu lengkungannya:
        \[
        L = \int_0^{2\pi} 2R\sin\frac t2\,\dd t = 8R
        \]
        (yaitu teorema Wren, 1658). Empat garis tengah roda, dan tanpa
        $\pi$ di mana pun.
  \item Hitunglah \omterm{def:b1:curves:arclength}{panjang busur} dari \omterm{ex:b1:curves:astroid}{titik runcingnya}: $s(t) =
        4R\bigl(1 - \cos\frac t2\bigr)$, lalu periksalah $s(2\pi) =
        8R$.
  \item Kini ukurlah busurnya dari \emph{puncaknya} $t = \pi$:
        $\sigma(t) = \abs{4R\cos\frac t2}$. Buktikan hubungan
        intrinsiknya
        \[
        \sigma^2 = 8R\,\bigl(2R - y\bigr) :
        \]
        bahwa kuadrat jarak busur dari puncaknya sebanding dengan
        penurunan tingginya di bawah puncaknya.
  \item Periksa kewarasannya: pulihkanlah dari pertanyaan 8 bahwa setengah lengkungan
        dari puncak ke \omterm{ex:b1:curves:astroid}{titik runcingnya} berpanjang $4R$, lalu bandingkanlah dengan
        perhitungan astroid \cref{ex:b1:curves:circlelength}
        --- karena kedua kurvanya berpanjang aljabar dan \omterm{def:b1:vspaces:free}{bebas} $\pi$;
        lalu jelaskanlah apa yang memungkinkan ini walaupun keduanya
        dibangun dari lingkaran. \emph{(Lihatlah bentuk
        $\norm{f'}$.)}
\end{enumerate}

\medskip
\noindent\textbf{Bagian III --- Luas Roberval: $3\pi R^2$.}
\begin{enumerate}[resume]
  \item Benarkanlah bahwa luas antara satu lengkungan dan tanahnya adalah
        $A = \int_0^{2\pi} y(t)\,x'(t)\,\dd t$ \emph{(yaitu
        penyulihan $x = x(t)$ pada $\int y\,\dd x$,
        \cref{thm:b1:integration:parts}; karena $x$ naik)}.
  \item Hitunglah
        \[
        A = R^2\int_0^{2\pi}(1 - \cos t)^2\,\dd t = 3\pi R^2 :
        \]
        yaitu tepat \emph{tiga} kali luas rodanya ---
        yakni nisbah yang telah ditebak Galileo lewat penimbangan.
  \item Metode yang sama bagi astroid $(\cos^3t, \sin^3t)$: tunjukkanlah
        bahwa luas yang dilingkupinya adalah $\frac{3\pi}8$
        \emph{(linearkanlah $\sin^4 t\cos^2 t$; karena hanya suku
        tetapnya yang bertahan atas satu periode penuh)}.
  \item Periksalah kewarasan rumus pertanyaan 10 pada setengah lingkaran
        satuan atas $(\cos t, \sin t)$, dengan $t$ dari $\pi$ ke $0$:
        apakah ia mengembalikan $\frac\pi2$?
\end{enumerate}

\medskip
\noindent\textbf{Bagian IV --- Tautokron Huygens.} Baliklah
lengkungannya: sebutir manik tanpa gesekan meluncur, di bawah gravitasi $g$, di dalam
mangkuk sikloidal
\[
x(t) = R(t + \sin t), \qquad y(t) = R(1 - \cos t)
\qquad (t \in \intcc{-\pi}{\pi}),
\]
yang titik terendahnya adalah titik asalnya ($y$ diukur ke atas).
\begin{enumerate}[resume]
  \item Hitunglah lajunya $\norm{f'(t)} = 2R\cos\frac t2$,
        \omterm{def:b1:curves:arclength}{panjang busur} dari dasarnya $s(t) = 4R\sin\frac t2$, lalu
        buktikanlah kesamaan kuncinya
        \[
        y = \frac{s^2}{8R} .
        \]
  \item Manik yang dilepas diam dari titik berparameter
        $t_0 > 0$ menuruti kekekalan energi: bahwa jika $s(\tau)$ menyatakan
        kedudukan busurnya pada waktu $\tau$, maka $\frac12\bigl(
        \frac{\dd s}{\dd\tau}\bigr)^2 + g\,y = g\,y_0$. Tulislah ulang
        ini, dengan memakai pertanyaan 14, sebagai
        \[
        \Bigl(\frac{\dd s}{\dd\tau}\Bigr)^{\!2}
        = \frac{g}{4R}\,\bigl(s_0^2 - s^2\bigr),
        \qquad s_0 = s(t_0).
        \]
  \item Turunkanlah terhadap $\tau$ lalu perolehlah
        pengayun harmonisnya
        \[
        \frac{\dd^2 s}{\dd\tau^2} = -\frac{g}{4R}\,s ;
        \]
        pecahkanlah ia dengan \cref{thm:b1:diffeq:homogeneous2} dan
        syarat awalnya: $s(\tau) = s_0\cos(\omega\tau)$,
        $\omega = \sqrt{g/4R}$.
  \item Simpulkanlah \emph{sifat tautokronnya} (Huygens, 1659):
        bahwa waktu untuk mencapai dasarnya,
        \[
        T_{\downarrow} = \frac{\pi}{2\omega}
        = \pi\sqrt{\frac Rg}\,,
        \]
        tak bergantung pada titik pelepasannya --- karena manik yang dilepas
        bersama dari sebarang dua ketinggian mangkuknya tiba
        bersama.
  \item Periksalah lewat penyulihan bahwa $s(\tau) =
        s_0\cos\omega\tau$ memenuhi persamaan energi berorde
        pertama pada pertanyaan 15 dengan persis (bukan hanya yang
        diturunkan), lalu jelaskanlah dalam satu kalimat mengapa
        bandul lingkaran hanya \emph{secara hampiran}
        isokron sedangkan sikloidnya persis demikian.
  \item Secara numeris: jari-jari $R$ yang mana yang membuat waktu turunnya
        tepat satu detik ($g = 9.81$)? Perhatikanlah betapa dekat
        jawabannya dengan satu meter, dan bagaimana periode ayunan
        penuhnya $2\pi\sqrt{4R/g}$ dibandingkan dengan rumus bandul
        bersudut kecil $2\pi\sqrt{\ell/g}$ untuk $\ell = 4R$.
\end{enumerate}

\medskip
\noindent\textbf{Bagian V --- Dividen, dan sintesisnya.}
\begin{enumerate}[resume]
  \item Terapkanlah kaidah singgung pertanyaan 3 di $t =
        \frac\pi2$ (ambillah $R = 1$): hitunglah $M$, $T$,
        arah $MT$, lalu periksalah ia terhadap
        $f'(\frac\pi2)$.
  \item (Trokoid) Tandailah sebagai gantinya sebuah titik berjarak $d$ dari
        pusatnya ($d \neq R$): maka kurvanya adalah $x = Rt -
        d\sin t$, $y = R - d\cos t$. Tunjukkan bahwa untuk $d < R$
        kurvanya reguler di mana-mana dan tak memperlihatkan perilaku
        singular apa pun ($x' > 0$: jadi ia maju), sedangkan untuk $d
        > R$ absisnya $x'$ berganti tanda dan kurvanya membuat
        simpul --- yakni roda kereta api bertepi yang titik peleknya
        bepergian \emph{mundur}.
  \item (Cuplikan brakistokron) Dari \omterm{ex:b1:curves:astroid}{titik runcing} $(\pi R, 2R)$ pada
        mangkuknya ke dasarnya, bandingkanlah waktu turun sikloidnya
        $\pi\sqrt{R/g}$ dengan waktu sepanjang peluncur lurus
        yang menghubungkan titik yang sama \emph{(dengan percepatan tetap
        $g\sin\alpha$ sepanjang tali busurnya)}: tunjukkanlah bahwa tali busurnya memakan
        $\sqrt{\pi^2 + 4}\,\sqrt{R/g} \approx 3.72\sqrt{R/g}$.
        Jadi kurvanya mengalahkan garisnya --- dan ia sesungguhnya yang tercepat
        di antara semua kurva, yaitu hasil kalkulus variasi.
  \item Tunjukkan bahwa pada lengkungannya ($0 < t < 2\pi$),
        \[
        \frac{\dd y}{\dd x} = \cot\frac t2,
        \qquad
        \frac{\dd^2 y}{\dd x^2} =
        -\frac{1}{4R\sin^4\frac t2} < 0 :
        \]
        bahwa lengkungannya cekung, dengan garis singgung tegak tepat pada
        \omterm{ex:b1:curves:astroid}{titik runcingnya}.
  \item Pulihkanlah garis singgung tegak \omterm{ex:b1:curves:astroid}{titik runcing} pada
        \cref{exo:b1:curves:2} secara geometri: hitunglah arah
        limit tali busur $MT$ pertanyaan 3 ketika $t \to
        0^{+}$, tanpa penjabaran sama sekali.
  \item Sintesis, dalam empat kalimat: ketiga teorema bernama yang mana
        yang dibuktikan soal ini (beserta bilangannya $8R$, $3\pi
        R^2$, $\pi\sqrt{R/g}$ dan penulisnya); perangkat perhitungan
        tunggal yang mana (yaitu pemfaktoran $\sin\frac t2$, $\cos\frac t2$)
        yang menggerakkan seluruh Bagian I, II dan IV; bagaimanakah
        hubungan intrinsik $y = s^2/8R$ mengubah geometri
        menjadi persamaan diferensial \omterm{def:b1:linmaps:def}{linear}; dan bab yang mana
        pada buku ini yang disandari tiap Bagiannya.

\end{enumerate}
\end{problem}
